% ══════════════════════════════════════════════════════════════════════════
% NEW REMARKS — insert after the Cartesian parametric representation
% (equations x(p), y(p)) in chapters/04_parametrisation.tex,
% before \section{Derivatives}
% ══════════════════════════════════════════════════════════════════════════

\begin{remark}[The angular factor as a Dirichlet $L$-value:
Eisenstein integers]
\label{rem:eisenstein_L}
The angular factor $\pi/\sqrt{3}$ of the $p$-parametrisation admits an
exact expression through a Dirichlet $L$-function. Let $\chi_{-3}$ be
the nontrivial character modulo~$3$, i.e.\ $\chi_{-3}(n) = +1$ for
$n \equiv 1 \pmod 3$, $-1$ for $n \equiv 2 \pmod 3$, and $0$ for
$3 \mid n$. Its $L$-series at $s=1$,
\begin{equation}\label{eq:L_chi3}
  L(1,\chi_{-3})
  = 1 - \frac12 + \frac14 - \frac15 + \frac17 - \frac18 + \cdots
  = \frac{\pi}{3\sqrt{3}},
\end{equation}
is a classical evaluation; by Dirichlet's class number formula it reads
$L(1,\chi_{-3}) = 2\pi h / (w\sqrt{|d|})$ with class number $h = 1$,
discriminant $d = -3$ and $w = 6$ units for the field
$\mathbb{Q}(\sqrt{-3})$~\cite{IrelandRosen1990}. Hence
\begin{equation}\label{eq:angular_L}
  \frac{\pi}{\sqrt{3}} = 3\,L(1,\chi_{-3}) = \frac{w}{2}\,L(1,\chi_{-3}),
  \qquad
  \theta(p) = \frac{w}{2}\,L(1,\chi_{-3})\,\sqrt{p}.
\end{equation}

The ring of integers of $\mathbb{Q}(\sqrt{-3})$ is the ring of
Eisenstein integers $\mathbb{Z}[\omega]$, $\omega = e^{2\pi i/3}$ ---
the hexagonal lattice, whose unit group has order exactly~$6$. Its
arithmetic is organised by the same congruence classes as the present
work: a rational prime $p > 3$ splits in $\mathbb{Z}[\omega]$ if
$p \equiv 1 \pmod 3$ and remains inert if $p \equiv 2 \pmod 3$, while
$p = 3$ ramifies. For odd $n$, the classes $n \equiv 1, 5 \pmod 6$ on
which $k(n)$ is integer-valued are precisely the classes
$n \equiv 1, 2 \pmod 3$, so the two strands of the twin structure
correspond to the two splitting types. The associated Dedekind zeta
function factorises as
$\zeta_{\mathbb{Q}(\sqrt{-3})}(s) = \zeta(s)\,L(s,\chi_{-3})$.

It should be stated clearly what is established and what is not: the
identities \eqref{eq:L_chi3}--\eqref{eq:angular_L} are exact, and the
correspondence between the residue classes and the splitting law in
$\mathbb{Z}[\omega]$ is classical algebraic number theory. That the
angular factor of the spiral is forced to be this particular $L$-value
--- rather than merely equal to it through the shared origin of both
quantities in the modulus~$6$ --- is a conjecture of the present work
and remains open.
\end{remark}

\begin{remark}[The angular factor and the reflection formula of the
Gamma function]
\label{rem:gamma_reflection}
Euler's reflection formula
$\Gamma(z)\Gamma(1-z) = \pi/\sin(\pi z)$ yields at $z = \tfrac13$
\begin{equation}\label{eq:gamma_reflection}
  \Gamma\!\left(\tfrac13\right)\Gamma\!\left(\tfrac23\right)
  = \frac{\pi}{\sin(\pi/3)}
  = \frac{2\pi}{\sqrt{3}},
  \qquad\text{hence}\qquad
  \frac{\pi}{\sqrt{3}}
  = \frac{1}{2}\,\Gamma\!\left(\tfrac13\right)\Gamma\!\left(\tfrac23\right).
\end{equation}
The arguments $\tfrac13$ and $\tfrac23$ are the two residue classes
modulo~$3$ --- normalised to the unit interval --- that carry the twin
structure of $k(n)$; the angular factor is one half of the product of
the Gamma values at exactly these two points.

The appearance of $\sqrt{3}$ in \eqref{eq:gamma_reflection} and in the
$L$-value of Remark~\ref{rem:eisenstein_L} is not independent: the
Chowla--Selberg formula~\cite{ChowlaSelberg1967} expresses the periods
attached to the imaginary quadratic field $\mathbb{Q}(\sqrt{-3})$ in
terms of $\Gamma(\tfrac13)$, so that the $L$-value and the
Gamma product are two evaluations rooted in the same field. Within the
present work these identities are recorded as structure shared with
the parametrisation constants; no derivation of the spiral from either
side is claimed.
\end{remark}
